\subsection{P001 --- Spatial-temporal CO2 concentration prediction in education facilities with dynamic layouts using BIM and multi-modal graph neural networks}
\label{paper:P001}
\textbf{Authors:} Hao Liu, Kam-Ming Mark Tam, Cong Huang, Chun Yin Lai, Eric Schuldenfrei\par
\textbf{Major Category:} Building Environment and IEQ\par
\textbf{Secondary Category:} AI and Machine Learning for Buildings, BIM and Semantic Information\par
\textbf{Keywords:} Spatial-temporal prediction, CO2 concentration, Education facility, Space reconfigurability, BIM, Graph neural networks, Multi-modal deep learning\par
\paragraph{主な主張}
BIM由来の時変レイアウトとVoronoiベースのCO2センサ隣接グラフをmulti-modal deep learningで統合し, 2時間先予測でMAE 11.25 ppmを達成し, 比較対象中で次点のTGCNの15.10 ppmより約25\%低いMAEを示した. \cite{Liu:2026DynamicCO2}
\paragraph{新規性}
静的レイアウトを前提とする既往のGNN型CO2予測に対し, BIMから時変の間仕切り構成をテンソル化し, センサグラフ文脈と空間構成文脈を二枝のMMDLで統合した点に新規性がある.
\paragraph{位置付け}
教育施設のIAQ予測を主目的としつつ, BIM幾何情報を予測モデルの入力文脈として利用するため, Building Environment and IEQを主分類, AI and Machine Learning for BuildingsおよびBIM and Semantic Informationを副分類として位置付ける.
